\chapter*{Remerciements}
\addcontentsline{toc}{chapter}{Remerciements}

Avant de présenter ce travail, je tiens à exprimer ma sincère gratitude envers toutes les personnes qui m'ont accompagné et soutenu tout au long de ce stage.

Je remercie tout d'abord \encadreurAcademique, mon encadreur académique à \ecole, pour sa disponibilité, ses conseils avisés et le suivi attentif qu'elle a assuré tout au long de ce projet. Ses remarques m'ont permis de prendre du recul sur mon travail et d'en améliorer la rigueur.

Je souhaite également adresser mes remerciements à \maitreDeStage, \maitreDeStageTitre, qui m'a accueilli au sein de son équipe et m'a fait confiance dès les premières semaines du stage. Son expérience et sa pédagogie ont grandement facilité ma compréhension des enjeux réels d'un projet informatique au service d'une administration publique.

Mes remerciements vont aussi à l'ensemble du personnel du \structureAccueil{} qui, par sa disponibilité et sa bienveillance, a rendu cette expérience aussi formatrice qu'agréable. Travailler à leurs côtés m'a permis de mesurer concrètement l'utilité du numérique dans la modernisation des services publics guinéens.

Je n'oublie pas le corps professoral de \ecole, dont la formation m'a donné les bases nécessaires pour mener ce projet à bien, ainsi que mes camarades de promotion pour leur soutien et les échanges enrichissants que nous avons eus durant ces années.

Enfin, je remercie ma famille et mes proches pour leurs encouragements constants, qui ont été un véritable moteur tout au long de mon parcours académique.
